% =====================================================================
% Figure contract
% 核心结论：Web、Mobile 与 Desktop OS 共享 GUI/CLI/API/MCP 动作通道，
%           Hybrid routing 将“选哪条通道”提升为独立决策问题。
% 图原型：原创三栏映射；平台 → 动作通道 → Hybrid 路由。
% 证据层级：hero = Hybrid 路由；支撑 = 3 个平台、4 个通道；收束 = 底部证据注释。
% 能不能更少？：不能；删任一平台或通道都会破坏指定 taxonomy，删注释会丢失 §4.8 的证据边界。
%
% Step ① 设计文档（Form B / Narrative，中等密度三栏图）
% 0 参考：沿用 cua-survey-fig1-overview.tex 的 XeLaTeX、低饱和 palette、圆角卡片与字号节奏。
% 1 领域：Computer-Use Agents 中文综述；模型名、接口名保持英文。
% 2 格式：standalone TikZ；fontspec + xeCJK；全图 sans-serif。
% 3 画布：W_ink <= 20cm；W_print = 19cm；目标印刷最小 CJK 字号约 7.5pt。
% 4 模块：左栏 3 张同色平台卡；中栏 4 张异色动作卡；右栏 1 张高权重 Hybrid 卡；底部整宽注释条。
% 5 连线：左侧 3→4 与右侧 4→1 均为 trunk + vertical spine + horizontal stubs；无对角线。
% 6 空间：两条 bus 各占栏间走廊；spine 与 stub 共享精确 named coordinate，接头零间隙。
% 7 强调：Hybrid 卡面积最大、唯一 1.5pt accent 描边并使用独有强调底色；其余描边 0.7pt。
% 8 密度：中等档；纯 taxonomy/route 结构，无可诚实嵌入的数值图表。
% 9 Narrative：中栏四卡构成垂直主轴；左栏围绕中栏中心放三卡；右侧 hero 对齐中栏中心；
%   标题与栏头占顶部一行，证据条紧贴主图底部，避免大块空白。
% 10 可视化决策：无数值，不嵌入 chart/heatmap；用双侧 bus 拓扑与 hero 层级作为记忆点。
%
% Pre-flight：P1 设计文档已写；P2 节点由 positioning/calc 相对构造；P3 所有连接均走 anchor；
% P4 N/A；P5 两条独立 spine 已分配；P6 annotation 统一全宽；P7 已读 lessons 与三栏专项规则。
% =====================================================================

\documentclass[border=8pt]{standalone}
\usepackage{fontspec}
\usepackage{xeCJK}
\usepackage{xcolor}
\usepackage{tikz}
\setCJKmainfont{Heiti SC}
\setCJKsansfont{Heiti SC}
\usetikzlibrary{arrows.meta,bending,positioning,calc}

% Low-saturation palette shared with the CUA survey figure batch.
\definecolor{ink}{HTML}{27364A}
\definecolor{muted}{HTML}{748092}
\definecolor{mainLine}{HTML}{4F667C}
\definecolor{blueLine}{HTML}{6D879E}
\definecolor{blueFill}{HTML}{E7EEF4}
\definecolor{purpleLine}{HTML}{817594}
\definecolor{purpleFill}{HTML}{EEEAF2}
\definecolor{goldLine}{HTML}{9B875B}
\definecolor{goldFill}{HTML}{F2EEE3}
\definecolor{tealLine}{HTML}{648D8A}
\definecolor{tealFill}{HTML}{E5F0EE}
\definecolor{roseLine}{HTML}{9A777B}
\definecolor{roseFill}{HTML}{F3E9EA}
\definecolor{noteFill}{HTML}{F1F3F5}
\definecolor{noteLine}{HTML}{A7AFB8}

\tikzset{
  every node/.style={
    font=\sffamily\fontsize{8.4}{10.2}\selectfont,
    text=ink
  },
  platform/.style={
    rectangle,
    rounded corners=4pt,
    draw=blueLine,
    fill=blueFill,
    line width=0.7pt,
    minimum width=2.35cm,
    minimum height=0.84cm,
    align=center,
    inner sep=5pt,
    font=\sffamily\fontsize{8.8}{10.4}\selectfont\bfseries
  },
  actioncard/.style={
    rectangle,
    rounded corners=5pt,
    line width=0.7pt,
    minimum width=6.75cm,
    minimum height=1.20cm,
    text width=6.35cm,
    align=left,
    inner xsep=8pt,
    inner ysep=5pt
  },
  hybridcard/.style={
    rectangle,
    rounded corners=7pt,
    draw=purpleLine,
    fill=purpleFill,
    line width=1.5pt,
    minimum width=5.05cm,
    minimum height=2.62cm,
    text width=4.48cm,
    align=center,
    inner sep=10pt
  },
  annotationbar/.style={
    rectangle,
    rounded corners=4pt,
    draw=noteLine,
    fill=noteFill,
    line width=0.6pt,
    minimum width=18.75cm,
    minimum height=1.02cm,
    text width=18.05cm,
    align=center,
    inner sep=6pt,
    font=\sffamily\fontsize{7.7}{9.5}\selectfont
  },
  % Zero-gap bus stub: no shortening at the spine-side start point.
  fan_stub/.style={
    -{Stealth[length=5pt 1.5,width'=0pt 0.6,sep=0pt 0.5]},
    line width=1.0pt,
    shorten >=2pt,
    shorten <=0pt,
    color=mainLine
  },
  busline/.style={
    line width=1.0pt,
    color=mainLine
  },
  rightbusline/.style={
    line width=1.05pt,
    color=purpleLine
  },
  colheader/.style={
    font=\sffamily\fontsize{8.8}{10.5}\selectfont\bfseries,
    text=mainLine
  },
  figtitle/.style={
    font=\sffamily\fontsize{11.0}{13.2}\selectfont\bfseries,
    text=ink
  }
}

\begin{document}
\sffamily
\begin{tikzpicture}

% Middle column is the construction anchor.
\node[actioncard,fill=blueFill,draw=blueLine] (gui)
  {{\fontsize{9.1}{10.7}\selectfont\bfseries GUI}\\[0pt]
   {\color{muted}\fontsize{7.5}{9.1}\selectfont 覆盖面最广;需 screenshot-action 视觉循环,较慢}};
\node[actioncard,below=0.24cm of gui,fill=tealFill,draw=tealLine] (cli)
  {{\fontsize{9.1}{10.7}\selectfont\bfseries CLI}\\[0pt]
   {\color{muted}\fontsize{7.5}{9.1}\selectfont 快,无 grounding 问题;仅覆盖 CLI 可达任务}};
\node[actioncard,below=0.24cm of cli,fill=goldFill,draw=goldLine] (api)
  {{\fontsize{9.1}{10.7}\selectfont\bfseries API}\\[0pt]
   {\color{muted}\fontsize{7.5}{9.1}\selectfont 稳定;上限取决于可用接口注入}};
\node[actioncard,below=0.24cm of api,fill=roseFill,draw=roseLine] (mcp)
  {{\fontsize{9.1}{10.7}\selectfont\bfseries MCP}\\[0pt]
   {\color{muted}\fontsize{7.5}{9.1}\selectfont 标准化协议;GUI 场景研究仍稀薄(§4.9)}};

\coordinate (actionCenter) at ($(cli)!0.5!(api)$);

% Left platform column, centered against the four-channel stack.
\node[platform,left=4.30cm of actionCenter] (mobile) {Mobile};
\node[platform,above=0.46cm of mobile] (web) {Web};
\node[platform,below=0.46cm of mobile] (desktop) {Desktop OS};

% Right hero card: area, fill, weight, and border all establish hierarchy.
\node[hybridcard,right=4.55cm of actionCenter] (hybrid)
  {{\fontsize{10.1}{12.0}\selectfont\bfseries Hybrid 路由(§4.8)}\\[5pt]
   {\fontsize{8.8}{10.9}\selectfont — 通道选择本身成为决策问题}};

% Column labels and overall title.
\node[colheader,above=0.38cm of web] (platformHeader) {平台};
\node[colheader,above=0.38cm of gui] (channelHeader) {动作通道};
\coordinate (titleCenter) at ($(web.west)!0.5!(hybrid.east)$);
\node[figtitle,anchor=south] at ([yshift=1.08cm]titleCenter |- gui.north)
  {平台 × 动作通道 × Hybrid 路由(§4.5–4.8)};

% Left 3→4 bus: three source trunks feed one continuous spine; four
% zero-gap stubs leave the spine orthogonally and terminate at card anchors.
\coordinate (leftSpineCenter) at ($(mobile.east)!0.5!(gui.west)$);
\coordinate (leftSpineTop) at (leftSpineCenter |- gui.west);
\coordinate (leftSpineBottom) at (leftSpineCenter |- mcp.west);
\draw[busline] (leftSpineTop) -- (leftSpineBottom);
\draw[busline] (web.east) -- (leftSpineCenter |- web.east);
\draw[busline] (mobile.east) -- (leftSpineCenter |- mobile.east);
\draw[busline] (desktop.east) -- (leftSpineCenter |- desktop.east);
\draw[fan_stub] (leftSpineCenter |- gui.west) -- (gui.west);
\draw[fan_stub] (leftSpineCenter |- cli.west) -- (cli.west);
\draw[fan_stub] (leftSpineCenter |- api.west) -- (api.west);
\draw[fan_stub] (leftSpineCenter |- mcp.west) -- (mcp.west);

% Right 4→1 bus: source stubs converge on one spine; one emphasized trunk
% carries the merged route into the Hybrid card. All joins share coordinates.
\coordinate (rightSpineX) at ($(gui.east)!0.5!(hybrid.west)$);
\coordinate (rightSpineCenter) at (rightSpineX |- actionCenter);
\coordinate (rightSpineTop) at (rightSpineCenter |- gui.east);
\coordinate (rightSpineBottom) at (rightSpineCenter |- mcp.east);
\draw[rightbusline] (rightSpineTop) -- (rightSpineBottom);
\draw[rightbusline] (gui.east) -- (rightSpineCenter |- gui.east);
\draw[rightbusline] (cli.east) -- (rightSpineCenter |- cli.east);
\draw[rightbusline] (api.east) -- (rightSpineCenter |- api.east);
\draw[rightbusline] (mcp.east) -- (rightSpineCenter |- mcp.east);
\draw[fan_stub,line width=1.35pt,color=purpleLine]
  (rightSpineCenter) -- (hybrid.west);

% Full-width evidence annotation.
\node[annotationbar,below=0.58cm of mcp.south] (annotation)
  {路由质量的瓶颈在负责分派的 Orchestrator 推理能力,而非执行模块本身(CoAct-1 backbone ablation);证据图见本节末};

\end{tikzpicture}
\end{document}
